\chapter{Literature Review}
\label{ch:literature}

\section{Scope and search procedure}

The review started from the 24 papers collected in the first phase of the project and was extended by keyword search and by following citations in both directions on arXiv, the ACL Anthology, OpenReview, Semantic Scholar and the proceedings sites of the main venues, up to October 2026. A paper was kept only if it does at least one of four things: it learns evidence selection with reinforcement learning in or near tables; it defines a table--text or multi-table environment or pipeline whose selection step a learned policy could replace; it supplies part of the sequential-decision machinery we need (scoring a variable set of candidates, an explicit stopping action); or it fixes an optimisation rule that a correct implementation has to respect. The title, author list, venue and identifier of every retained paper were checked on its arXiv, Anthology, OpenReview or publisher page. More than one hundred candidates were checked in this way. Work that turned out to be adjacent but not relevant was dropped, for example prompting-only table agents, text-to-SQL reinforcement learning that always receives the full database schema, and multimodal table models.

The retained work falls into four groups, discussed in Sections~\ref{sec:lit-graph}--\ref{sec:lit-search}: reinforcement learning over graphs, reinforcement learning for table reasoning, table--text and multi-table question answering systems, and RL-trained evidence selection for text. Section~\ref{sec:lit-data} covers benchmarks, and Table~\ref{tab:landscape} at the end of the chapter places the closest papers side by side.

\section{Reinforcement learning for multi-hop reasoning over graphs}
\label{sec:lit-graph}

The idea of treating multi-hop reasoning as a walk whose next step is chosen by a policy comes from knowledge-graph completion. DeepPath \citep{xiong2017deeppath} cast relation-path finding as an MDP solved with REINFORCE, with a reward combining accuracy, path efficiency (the inverse of path length) and diversity. Two of its design choices recur in our work. First, the policy was pre-trained by imitating breadth-first-search paths, and the authors report that without this the model ``will show very poor convergence properties''. Second, the efficiency term is an early form of a per-hop cost. DeepPath, however, needs both the source and the target entity, so it solves link prediction rather than question answering.

MINERVA \citep{das2018minerva} removed the known target. The agent starts at the query entity, an LSTM summarises the walk so far, and every outgoing edge, represented by the concatenation of its relation and target-entity embeddings, is scored against a projection of the history and the query relation. The scores are normalised with a softmax over whatever edges happen to exist at the current node, which makes the policy an order-invariant pointer over a variable-size set; this is the template for our actor. MINERVA did not have an explicit stopping action: it added a self-loop and a fixed horizon, noting that a special STOP action was possible but unnecessary for their data. It was trained with REINFORCE, 20 rollouts per example and a moving-average baseline. Because entities are represented by learned identifier embeddings, test triples involving unseen entities had to be removed, which shrank the NELL-995 test set from 3{,}992 to 2{,}818 queries.

\citet{lin2018multihop} kept this MDP and added three devices: soft reward shaping, $R=R_b+(1-R_b)f$, where $f$ is a pre-trained embedding model that scores answers the graph does not contain; \emph{action dropout}, which randomly masks outgoing edges during sampling to fight spurious paths; and neighbour caps based on PageRank. Together they raised FB15k-237 MRR from MINERVA's 29.3 to 40.7. In their ablation, removing action dropout costs more (dev MRR 38.2 to 31.0) than removing shaping (38.2 to 35.1). Their warning about spurious paths, ``there are usually more spurious paths than correct ones, spurious paths are often found first'', applies directly to tables, where the same value can appear in many rows.

M-Walk \citep{shen2018mwalk} is the only classic walker with an explicit STOP action, scored by a separate network from a summary of the current neighbourhood, while neighbours are scored by an inner product; the policy and the Q-values share these scores. It was trained with Monte-Carlo tree search and Q-learning rather than policy gradients and reached 89.9 MAP on NELL-995 against MINERVA's 87.6. The authors note a hard ceiling that also applies to us: when the correct answer is not among the candidates the walker can reach, it always fails. SRN \citep{qiu2020srn} brought RL path search to natural-language KBQA with answer-only supervision and used potential-based reward shaping to deal with the delayed, sparse reward; only its abstract-level description could be verified.

Recent work trains large language models as graph agents. KG-R1 \citep{lin2025kgr1} uses a schema-agnostic action interface and GRPO, and rewards the agent for retrieving any gold entity; removing that reward costs 4.03 and 2.68 Hits@1 on WebQSP and CWQ. The same paper reports that replacing GRPO by PPO dropped WebQSP from 82.75 to 59.62 and CWQ from 65.31 to 35.66, with reward hacking in which the agent fabricated retrieved content. Graph-R1 \citep{luo2026graphr1} also found GRPO ahead of PPO and REINFORCE++. Search-on-Graph-R1 \citep{sun2026sogr1} shows how much a supervised warm start matters: CWQ reaches 88.7 with supervised fine-tuning followed by RL but only 63.0 with RL alone. GraphRAG-R1 \citep{yu2026graphragr1} introduces a cost-aware F1 reward that pays more for the same F1 when fewer retrievals are used. All of these PPO-versus-GRPO comparisons concern token-level LLM policies whose critic must read long text; none of them tests a compact pointer policy with a learned critic.

Taken together, the graph literature supplies the pointer head, the STOP action, shaping, action dropout and the warm-start recipe. What does not carry over is the representation. These walkers assume a fixed relation vocabulary and identifier embeddings for entities. In HybridQA every test table is new and every column header plays the role of a new relation, so candidates must be represented by their content, and the state must describe a growing \emph{set} of evidence rather than a single current node.

\section{Reinforcement learning for table reasoning}
\label{sec:lit-table}

\paragraph{Semantic parsing.} Before large language models, RL on tables meant learning programs from answer-only supervision. MAPO \citep{liang2018mapo} kept a memory buffer of high-reward programs and explored systematically, reaching 46.3\% on WikiTableQuestions with an ensemble of ten models (43.8\% for a single model). Its authors warn that ``high-reward but spurious programs will bias the gradient''. MeRL \citep{agarwal2019merl} meta-learned an auxiliary reward that discounts accidentally correct programs and improved the single model from 42.9\% to 44.1\%. These papers learned which program to run, not which evidence to read, but their concern with right answers obtained for the wrong reason is the same one that motivates our evidence-recall signal.

\paragraph{Selecting part of one table.} Three recent papers reward a model for where it looks inside a single table, and each makes the selection once. Learning to Reduce \citep{lee2024learningtoreduce} trains two FLAN-T5 models with PPO to output the relevant columns and then the relevant rows for a frozen GPT-4 reader. The reward is asymmetric (missing a necessary item costs much more than including an irrelevant one), and gold evidence is obtained by deleting rows and columns from WikiTableQuestions tables as long as the SQUALL SQL query still returns the answer. Applied zero-shot to HybridQA tables, the column selector reaches 87\% recall, against 70\% for a RoBERTa classifier and about 71.5\% for zero-shot GPT-4. The reward is computed on the selection, not on the downstream answer, and linked passages are ignored. The region-based Table-R1 \citep{wu2025tabler1region} makes the model emit row and column index sets inside its chain of thought and trains with a GRPO variant on a reward $\alpha_t\,\mathrm{IoU}+(1-\alpha_t)\,\mathrm{answer}$ whose region weight $\alpha_t$ decays during training. It reports an average gain of 14.36 points across models and benchmarks and 67.5\% fewer generated tokens than plain GRPO, and states that multi-table settings were not validated. Reasoning-Table \citep{lei2025reasoningtable} pays for cited cells through $R=R_{\mathrm{ans}}(1+\lambda_1R_{\mathrm{pos}})+\lambda_2R_{\mathrm{fmt}}$, so evidence counts only when the answer is right; on OTT-QA the gold table is supplied, so table choice is never learned.

\paragraph{Reasoning over whole tables.} Most of the 2025--2026 work on RL for tables does not select evidence at all. The inference-time-scaling Table-R1 \citep{yang-etal-2025-table} runs GRPO over fully serialised tables, and its Llama-3.1-8B model reaches 81.2 on WikiTableQuestions against 68.0 for GPT-4.1. TableMind \citep{jiang2026tablemind} trains a Python-tool agent with a GRPO variant and lists multi-table reasoning as future work. Formula-R1 \citep{cao2025formular1} is a rare PPO user, with a reward of 1 for a correct answer, 0.2 for an executable but wrong formula and 0 otherwise.

\paragraph{Selecting tables.} OpenTable-R1 \citep{qiu2025opentabler1} is the one paper that genuinely learns, with RL, which tables to fetch: a Qwen3-4B agent issues BM25 table searches and SQL calls and is trained with asynchronous GRPO. Exact match on open-domain WikiTables rises from 12.0 zero-shot to 72.1 after a supervised cold start and to 86.2 after RL. Each question has a single gold table and no joins are involved. In text-to-SQL, Reasoning-SQL \citep{pourreza2025reasoningsql} gives partial credit for the overlap between the tables and columns a query uses and the gold ones, which is the only RL credit for table choice we found there, and it is assigned once for the whole query. RE-Tab \citep{kwok2026retab} proposes a deterministic reward over the intermediate sub-tables produced by Chain-of-Table-style operations; only its abstract could be verified.

OpenTable-R1 rules out any claim that RL-based table selection is unexplored. The defensible claim is narrower: entity-level, multi-hop selection across rows, cells and linked text, with a learned decision to stop.

\section{Multi-hop table--text and multi-table question answering}
\label{sec:lit-hybrid}

\paragraph{Supervised pipelines on HybridQA.} HybridQA \citep{chen-etal-2020-hybridqa} pairs each question with a Wikipedia table whose cells are hyperlinked to passages, and most questions need one hop from a table row to a linked passage or back. The accompanying model, HYBRIDER, links question spans to cells, ranks the linked cells with a BERT classifier, hops to a neighbouring cell or passage, and finally reads. Its stages are trained independently on reasoning paths synthesised from string matches, which the authors call ``somewhat noisy'', and inference is greedy: ``we use greedy search at first two steps to remain only the highest probably cell''. On the dev set the stage accuracies are 87.4\% for linking, 87.9\% for ranking, 89.2\% for hopping and 62\% for reading, and the large model scores 44.0 EM / 50.7 F1 (43.8/50.6 on test). Later systems improved the components but kept the same pattern of weakly supervised, myopic selection. DocHopper \citep{sun2021iterative} attends iteratively over evidence (47.7/55.0 dev). MuGER$^2$ \citep{wang-etal-2022-muger2} compares a table score with a passage score (57.1/67.3 dev). MITQA \citep{kumar-etal-2023-multi} keeps a fixed top-5 rows, which by itself raises gold-row recall by 8--9\% over keeping one row (65.5/72.7 dev). TACR \citep{wu-etal-2023-tacr} selects cells and reports being ``relatively worse in row selection'' (64.5/71.6 dev). S$^3$HQA \citep{lei-etal-2023-s3hqa} uses rules keyed on the question type and loses 2.9 EM when that selector is removed (68.4/75.3 dev, the best fully supervised result we found). None of these systems is trained against the final answer.

\paragraph{Open-domain table--text QA.} OTT-QA \citep{chen2021open} removes the gold table and the hyperlinks at test time. Its authors abandoned iterative retrieval because ``early mistakes can propagate to later retrieval stages''. CORE \citep{ma2022open} and Chain-of-Skills \citep{ma2023chain} chain retrieved tables and passages under fixed budgets (100 tables, 200 rows, top-100 passages for Chain-of-Skills), reaching 63.2 F1 on dev with Chain-of-Skills.

\paragraph{LLM-based frameworks with explicit selection steps.} HPROPRO \citep{shi-etal-2024-exploring} answers HybridQA questions with program-based prompting (48.7/57.7 on the full test set). The Hybrid Graphs paper \citep{agarwal-etal-2025-hybrid}, whose method is called ODYSSEY, links question-relevant cells to named entities and their passages, keeps the highest-ranked entity whose embedding similarity exceeds a fixed threshold of 0.8, and prunes the graph by a breadth-first traversal to depth three. With GPT-4 it reaches 58.4 EM / 71.8 F1 on a 500-question sample of HybridQA dev, against 55.4 EM when the full table and passages are given, and 14 of its 41 analysed errors come from the selection module. N2N-GQA \citep{sharafath2026n2ngqa} prunes a similar graph with a centrality score and fixed node budgets (41.50/48.17 on a 200-question sample). Among single-table frameworks, Chain-of-Table \citep{wang2024chainoftable} lets the LLM choose the next table operation, including row and column selection, by a greedy search, and reaches 86.61 on TabFact and 67.31 on WikiTableQuestions with PaLM~2. TableRAG \citep{chen2024tablerag} retrieves schema entries and cell values by embedding similarity under a token budget, and its 49.2 versus 37.7 lead over row-and-column retrieval on ArcadeQA comes from what is selected. Dater \citep{ye2023dater} prompts the LLM to reduce a large table to a sub-table, and Binder \citep{cheng2023binder} leaves selection implicit in a generated program.

\paragraph{Multi-table retrieval.} MMQA \citep{wu2025mmqa} is the only multi-table, multi-hop benchmark with gold SQL and key annotations (2{,}591 two-table and 721 three-table questions built from Spider). Its retriever multiplies relevance by a column-overlap score and stops when the overlap is zero; the best model, o1-preview, reaches 50.78 EM on two-table and 48.28 on three-table questions. JAR \citep{chen2024jar} solves a mixed-integer program over relevance and joinability, and a greedy join-aware search matches it while running 4 to 400 times faster \citep{boutaleb2025exploring}, which suggests that a sequential formulation is not intrinsically weaker than global optimisation. MURRE \citep{zhang2025murre} retrieves one table per hop with beam search, lets an LLM rewrite the question to ``remove'' what is already covered and to output ``None'' when it should stop, caps the search at three hops, and leaves learning from downstream text-to-SQL feedback to future work. Every method degrades with depth: on a unified MMQA corpus, recall@2 drops from two-hop to three-hop questions from 55.90 to 44.70 for JAR and from 49.15 to 34.05 for MURRE \citep{ho2026pearl}. REaR \citep{agarwal2025rear} explicitly suggests reinforcement learning for choosing the expansion depth.

\section{RL-trained evidence selection and search agents}
\label{sec:lit-search}

\paragraph{Before large language models.} Supervised path retrieval was more successful than RL in text QA. \citet{asai2020learning} retrieve reasoning paths over the Wikipedia hyperlink graph with a recurrent retriever and learn when to stop by supervision: an end-of-evidence symbol is a negative candidate at every step before the gold path ends. \citet{das2019multistep} trained a query reformulator with RL against a fixed reader and found that a minibatch-average reward baseline ``significantly degrades the final performance'' because questions differ in difficulty, an argument for per-question or learned baselines.

\paragraph{The closest precedent.} \citet{jose2024qatexttablesdrl} are the only authors we found who train an RL controller for table--text QA. On OTT-QA, a PPO or DQN agent chooses at each step between \emph{retrieve texts}, \emph{retrieve tables} and \emph{generate answer}, for at most four steps, with a frozen retriever and a frozen Fusion-in-Encoder reader. The reward is 2 for an exact match, the token F1 for a partial match and $-0.5$ otherwise, minus 0.02 per retrieval. The best agent (PPO with a Transformer feature extractor over BM25) scored 14.81 EM / 19.03 F1, identical to the fixed ``always retrieve text'' sequence; in the authors' words, it was ``converging to the trivial solution of only picking the best action on average''. Chapter~\ref{ch:gaps} analyses why.

\paragraph{RL search agents.} The 2025--2026 search agents train a language model to interleave reasoning with retrieval. Search-R1 \citep{jin2025searchr1} reports that GRPO converges faster than PPO but suffers reward collapse after many steps, whereas PPO remains stable, and that masking retrieved tokens out of the loss helps consistently. A controlled follow-up study by the same group \citep{jin2025empirical} found that adding intermediate retrieval rewards ``does not significantly improve final performance for either PPO or GRPO'' and that larger weights hurt. The opposite evidence is equally specific. StepSearch \citep{wang2025stepsearch} trains with PPO and GAE on a per-round reward equal to the gain in coverage of the gold documents minus a redundancy penalty, plus answer F1, and gains 5.7 to 15.2 points over standard RL baselines on four multi-hop benchmarks. IGPO \citep{wang2026igpo} shows that outcome-only GRPO suffers ``advantage collapse'' when all rollouts of a question score the same, and its turn-level information-gain reward adds 16.6 points for a 3B model and 8.3 for a 7B model.

The design closest to ours in spirit is s3 \citep{jiang2025s3}. Its searcher generates queries, selects a subset of the retrieved documents and decides whether to continue; a \emph{frozen} generator then answers, and only the searcher is trained, with PPO. The reward, ``Gain Beyond RAG'', is the generator's accuracy with the selected evidence minus its accuracy with naive top-$k$ retrieval. s3 needed about 2.4k training examples and 20 PPO steps, against 170k examples and 2{,}100 steps for Search-R1. MMOA-RAG \citep{chen2025mmoarag} restricts a document selector's actions to document identifiers, which it says greatly reduces the exploration space, and follows the same supervised-then-PPO recipe that we use. Search Wisely \citep{wu2025searchwisely} measures how often agents search too little or too much and reduces under-searching from 42.04\% to 34.71\% of cases with its $\beta$-GRPO variant, which shows that the stopping decision is a failure mode in its own right.

\section{Benchmarks}
\label{sec:lit-data}

HybridQA has 62{,}682 training, 3{,}466 dev and 3{,}463 test questions over roughly 13{,}000 Wikipedia tables; an average table has about 15.7 rows, 4.4 columns and 44 hyperlinked passages, and 56.4\% of answers lie in passages, 42.3\% in cells and 1.1\% must be computed \citep{chen-etal-2020-hybridqa}. Gold evidence exists only in the automatically traced files, which record every cell or linked passage that contains the answer string. The OTT-QA authors estimate that this tracing is about 15\% wrong \citep{chen2021open}, and on the dev set only 1{,}890 of the 3{,}466 questions (54.5\%) have exactly one traced location. Test answers are hidden behind a leaderboard. Table~\ref{tab:hybridqa-published} lists published results; they come from fine-tuned or LLM readers and are not directly comparable to the frozen extractive reader used here, but they place the numbers of Chapter~\ref{ch:results} in context.

OTT-QA needs open retrieval over more than 400{,}000 tables and five million passages, which is beyond our compute budget. MMQA's official data link has expired; the only copy we found is an unlicensed third-party mirror with a single test split and tables averaging between 1{,}369 and 1{,}833 rows, so it can serve at most as an optional transfer experiment. WikiTableQuestions \citep{pasupat-liang-2015-compositional} provides answers but no evidence annotation, TAT-QA \citep{zhu-etal-2021-tat} is dominated by arithmetic over small financial tables, and MultiTabQA \citep{pal-etal-2023-multitabqa} has table-valued answers that do not fit a span reader. HybridQA is therefore our primary benchmark.

\begin{table}[t]
\centering
\small
\caption{Published HybridQA results (EM / F1). LLM results marked $^\dagger$ use 200- or 500-question samples of the dev set. All systems use fine-tuned or large language model readers.}
\label{tab:hybridqa-published}
\begin{tabular}{@{}llcc@{}}
\toprule
System & Selection mechanism & Dev & Test \\
\midrule
HYBRIDER (BERT-large) \citep{chen-etal-2020-hybridqa} & greedy rank-then-hop & 44.0 / 50.7 & 43.8 / 50.6 \\
DocHopper \citep{sun2021iterative} & iterative attention & 47.7 / 55.0 & 46.3 / 53.3 \\
MuGER$^2$ (large) \citep{wang-etal-2022-muger2} & score comparison rule & 57.1 / 67.3 & 56.3 / 66.2 \\
MITQA \citep{kumar-etal-2023-multi} & fixed top-5 rows & 65.5 / 72.7 & 64.3 / 71.9 \\
TACR \citep{wu-etal-2023-tacr} & cell classifier & 64.5 / 71.6 & 66.2 / 70.2 \\
S$^3$HQA \citep{lei-etal-2023-s3hqa} & rules by question type & 68.4 / 75.3 & 67.9 / 75.5 \\
HPROPRO, GPT-4 \citep{shi-etal-2024-exploring} & program prompting & 48.0 / 54.6$^\dagger$ & 48.7 / 57.7 \\
ODYSSEY, GPT-4 \citep{agarwal-etal-2025-hybrid} & threshold + BFS on graph & 58.4 / 71.8$^\dagger$ & -- \\
\bottomrule
\end{tabular}
\end{table}

\begin{table}[p]
\centering
\footnotesize
\caption{The closest prior work against the properties this project combines. ``Seq.'' means that evidence is selected one item per step with a stopping decision.}
\label{tab:landscape}
\setlength{\tabcolsep}{3.5pt}
\begin{tabularx}{\textwidth}{@{}>{\raggedright\arraybackslash}p{2.5cm}>{\raggedright\arraybackslash}X>{\raggedright\arraybackslash}X>{\raggedright\arraybackslash}p{1.6cm}>{\raggedright\arraybackslash}p{2.05cm}>{\raggedright\arraybackslash}p{1.8cm}>{\raggedright\arraybackslash}p{1.1cm}@{}}
\toprule
Work & Evidence & Action unit & Seq.\ + stop & Algorithm & Evidence reward & Frozen reader \\
\midrule
\citet{jose2024qatexttablesdrl} & tables + text (OTT-QA) & module: retrieve text / tables / answer & yes (answer) & PPO, DQN & none & yes \\
Learning to Reduce \citep{lee2024learningtoreduce} & one table & column set, then row set & no & PPO & asymmetric relevance & yes \\
Table-R1, region \citep{wu2025tabler1region} & one table & row and column indices in the reasoning chain & no & GRPO variant & region IoU, decaying & no \\
Reasoning-Table \citep{lei2025reasoningtable} & one table & cited cells & no & GRPO & cell Jaccard, answer-gated & no \\
OpenTable-R1 \citep{qiu2025opentabler1} & table corpus & search and SQL calls & multi-turn & GRPO & none & no \\
s3 \citep{jiang2025s3} & text & query, document subset, stop & yes & PPO & gain over naive RAG & yes \\
StepSearch \citep{wang2025stepsearch} & text & search round & implicit & PPO + GAE & gold-document gain & no \\
MINERVA \citep{das2018minerva} & knowledge graph & next edge & self-loop & REINFORCE & none & -- \\
M-Walk \citep{shen2018mwalk} & knowledge graph & next node or STOP & yes & MCTS + Q-learning & none & -- \\
HYBRIDER \citep{chen-etal-2020-hybridqa} & table + passages & cell, then cell or passage & greedy, 2 steps & none & -- & no \\
ODYSSEY \citep{agarwal-etal-2025-hybrid} & table--text graph & graph nodes & BFS, depth 3 & none & -- & yes (LLM) \\
MURRE \citep{zhang2025murre} & many tables & one table per hop & beam, LLM stop & none & -- & yes (LLM) \\
\midrule
\textbf{This project} & table + passages (extensible to joins) & one row or passage per step & yes, learned & PPO + GAE after BC & gold-path recall (shaping or bonus) & yes \\
\bottomrule
\end{tabularx}
\end{table}
